\documentclass[11pt]{article}
\usepackage[margin=1in]{geometry}
\usepackage{booktabs}
\usepackage{hyperref}
\usepackage{amsmath}
\usepackage{longtable}

\title{Executable Falsification for AI-Assisted Research:\\
Construction-Matched Controls, a Failure Taxonomy,\\
and a Research Ledger}
\author{Giovanni Pezzin \\ Physis Research Program}
\date{}

\begin{document}
\maketitle

\begin{abstract}
When AI agents design, implement, execute and interpret their own
experiments, the loop that produces results also produces the errors that
make those results wrong. We report twelve documented cases from a
single research programme in which a positive result was killed by a
control, and we extract from them a control taxonomy, a failure taxonomy,
and an executable workflow. Our central instrument is the
\emph{construction-matched null}: a control generated by a procedure that
preserves irrelevant aspects of the experimental construction while
selectively destroying or perturbing the mechanism claimed to produce the
effect. The distinction matters because our own first definition of it was
wrong, and we report that error and its correction. We show a case in
which only the null changed --- from random injection to nearest-cell
injection --- and the measured performance moved from AUC 0.882 to 0.506.
We then formalise the workflow that produced these catches: claims are
registered before running, with a pre-registered null; the null runs in
the same execution as the treatment; a verdict is assigned from the result;
revisions preserve the superseded claim; and failed branches are retained
rather than deleted. We are explicit about what this workflow does
\emph{not} do: it records protocol violations but cannot prevent them, and
we have no control group showing that it produces better research than
research without it.
\end{abstract}

\section{Introduction}

An AI agent that designs an experiment, writes the code, runs it, reads the
output and writes the conclusion concentrates four normally-separated
roles into one loop that executes in minutes. The efficiency is the point.
The hazard is that each of those four roles is a place where an error can
enter, and a single machine-mediated loop can generate and propagate such
errors far faster than a human-paced process.

This paper asks what actually goes wrong, and what instrumentation catches
it. It is deliberately not a success story: the most useful thing the
programme produced was a series of strong positive results that turned out
to be artifacts of their own controls.

We make three contributions, in three layers:

\begin{itemize}
\item \textbf{A failure catalog} (Section~\ref{sec:catalog}): twelve cases,
  each with the intended claim, the apparent positive, the failure mode,
  the control that caught it, the corrected result, the control's type, and
  when in the process it was detected.
\item \textbf{A control taxonomy} (Section~\ref{sec:taxonomy}): a precise
  definition of the construction-matched null, and an explicit account of
  what is \emph{not} one --- because stretching the term to cover every good
  research practice makes it useless.
\item \textbf{An executable workflow} (Section~\ref{sec:workflow}): the
  claim/experiment/verdict/revision structures, the invariant that no result
  may modify its own decision rule after observation, and the retention of
  failed branches, demonstrated on a real lineage.
\end{itemize}

\section{The null is the instrument}
\label{sec:instrument}

Consider the canonical experiment in our programme. A system claims to
detect misfiled ontology entries at AUC 0.880. Whether that number means
anything depends entirely on where the test errors were placed.

If errors are placed uniformly at random across cells, a large displacement
results, and any distance-based scorer detects them. If errors are placed
at the entry's nearest alternative cell --- a plausible mistake of the kind
a person actually makes --- the displacement is small and the scorer has
little to work with.

We ran both. AUC moved from 0.880 to 0.506. The system was identical. The
metric was identical. Only the construction of the null changed, and it
changed the conclusion from ``working detector'' to ``chance''.

\subsection{A definition we got wrong}

An earlier draft of this work defined a construction-matched null as one
that ``places errors where the system cannot detect them.'' The same draft
then stated correctly, in its principles section, that a null placing
errors where the system cannot detect them ``is not a control; it is a
tautology.'' Both sentences cannot hold.

The second is right. A null whose perturbations are undetectable by
construction guarantees the result and tests nothing. The correct
definition is:

\begin{quote}
A \textbf{construction-matched null} is a control generated by a procedure
that preserves irrelevant aspects of the experimental construction while
selectively destroying or perturbing the mechanism claimed to produce the
effect.
\end{quote}

The operative words are \emph{preserves} and \emph{selectively}. A valid
control must keep everything that is not the claim --- the data
distribution, the nuisance structure, the difficulty of the task --- and
vary only the mechanism. Applied to misfiling:

\begin{itemize}
\item \textbf{Bad control:} relocating an entry to a random cell. This
  fails to preserve difficulty (the displacement is unrealistic) and so the
  ``detection'' is of distance moved, not of misfiling.
\item \textbf{Better control:} relocating an entry to a semantically
  confusable alternative cell, preserving the displacement scale of real
  mistakes.
\end{itemize}

The control should kill the mechanism, not make the task impossible by
fiat. Getting this backwards is itself one of the failure modes we
catalog.

\section{Layer A: the failure catalog}
\label{sec:catalog}

We group the twelve cases by the phase in which the error entered, and
record for each the type of control that caught it and when detection
occurred. Table~\ref{tab:catalog} is the full catalog; the subsections give
the cases in enough detail to be re-used.

\begin{longtable}{@{}p{0.5cm}p{4.0cm}p{3.4cm}p{3.0cm}p{3.2cm}@{}}
\caption{The twelve cases. \textbf{CT} is control type; \textbf{Det} is
detection status.}\label{tab:catalog}\\
\toprule
\textbf{\#} & \textbf{Claim} & \textbf{Apparent positive} & \textbf{Control} & \textbf{Corrected / CT / Det} \\
\midrule
\endfirsthead
\toprule
\textbf{\#} & \textbf{Claim} & \textbf{Apparent positive} & \textbf{Control} & \textbf{Corrected / CT / Det} \\
\midrule
\endhead
1 & Geometry certifies misfiling (E14) & AUC 0.880 & nearest-cell injection & 0.506 \newline mechanism null \newline detected after \\
2 & Antonym rule detects contradiction (E83) & recall 0.996 & hand-written fixtures with no lexical antonym & 0.143 \newline adversarial fixture \newline detected after \\
3 & Structured rules precise on contradictions (E85) & precision 1.000 & fixtures built to fire each rule wrongly & 0.000 \newline adversarial fixture \newline detected after \\
4 & Five congruent domains aid abstraction (E19b) & C5 $-$ A1 $= +0.367$ & coupled-domain arm B5c & control invalid at 4B, valid at 9B \newline negative control \newline retrospectively \\
5 & Cline extraction is reliable (E74) & failure rate 0.005 & inspect shell-action output & 0.070 \newline measurement sanity check \newline detected during \\
6 & Corpus holds 402 sessions (E73) & 402 sessions & scan for byte-identical session ids & 188 duplicates dropped \newline replication check \newline detected during \\
7 & File revisions form a tight region (E78) & within-object distance 0.0000 & check embedding covers whole file & whole-revision chunking; 0.0118 \newline implementation invariant \newline detected during \\
8 & Coherence rises after perturbation (E76) & 4.0 $\rightarrow$ 5.0 & check which field the parser reads & parser read node count \newline measurement sanity check \newline detected during \\
9 & ANCHORVOC is not circular (E17) & tests 1--2 pass & three pre-registered tests, all required & test 3 fails; artifact \newline mechanism null \newline detected after \\
10 & Pooled baseline scores 0.400 (E87) & 0.400 & rerun baseline in the same script & 0.200 \newline in-run baseline \newline detected during \\
11 & Static rules find no structure (E81) & ratio 0.96 & decompose by cue arity & multi-word 15:1, single-word 0.92:1 \newline measurement sanity check \newline detected after \\
12 & The experiment is powered (E10) & pre-flight gate & power analysis & STOP returned, run proceeded \newline power gate \newline prevented output, not execution \\
\bottomrule
\end{longtable}

\subsection{Design-phase errors}

\textbf{Case 1 (E14).} The injection distribution is the null. Covered in
Section~\ref{sec:instrument}: $0.880 \rightarrow 0.506$ on a change of null
alone.

\textbf{Case 2 (E83).} A fixture generated by inserting a WordNet antonym
into a sentence produces contradictions that the antonym rule discovers by
construction. On the generated set the combined antonym+negation rules
reach recall 0.996 (F1 0.998). On hand-written fixtures carrying no lexical
antonym, recall falls to $0.143$ --- one contradiction of seven. The same
rule, two fixtures, two very different numbers: any figure from a
same-process fixture is an upper bound. We quote recall 0.996 rather than
the more flattering F1, because recall is the quantity that collapses.

\textbf{Case 3 (E85).} The same rule set reaching precision 1.000 on a
benign holdout collapses to precision 0.000 when given fifteen consistent
pairs shaped to trigger each rule. The stored artifact records five rule
firings --- quantifier 2, comparative 1, antonym 1, domain 1 --- every one
of them wrong. (The prose report of this experiment described four pairs
firing; the artifact records five firings and does not record a pair count,
so we cite the artifact and note the discrepancy rather than pick the
tidier number.) The failure has one shape: rules match class cues without
binding them to the pair's participants.

\textbf{Case 4 (E19b).} The control arm was invalid, and its validity
depended on the model. Arm B5c reached $+0.233$ at 9B, where it is a valid
control, but fired at 4B, 7B and 9B. Control validity is a property of the
stimulus and the model together, and must be checked per model.

\subsection{Implementation-phase errors}

\textbf{Case 5 (E74).} Reading only one key of Cline's shell results
yielded a failure rate of 0.005 over 8,301 actions; reading the other
result keys yielded 0.070. The pipeline reported an agent that never failed
because the parser never looked.

\textbf{Case 6 (E73).} Forty-two Claude Code sessions appeared on three
machines; 18 were byte-identical copies. The corpus count was right and the
session count was inflated.

\textbf{Case 7 (E78).} Embedding only the first 1,200 bytes of each
revision produced identical vectors for twelve revisions and a within-object
distance of exactly 0.0000. The ``tight region'' was a constant function.

\textbf{Case 8 (E76).} A parser selected the first number following the
word ``coherence'' and reported a rise from 4.0 to 5.0. The number was the
node count.

\subsection{Interpretation-phase errors}

\textbf{Case 9 (E17).} Three pre-registered confound tests disagreed: two
passed, one failed. The first verdict reported only the first test and
declared the confound resolved. Requiring all three to pass changed the
conclusion. A confound check that stops at the first reassuring answer is
not a check.

\textbf{Case 10 (E87).} A baseline of 0.400 cited from an earlier report
read 0.200 when rerun in the same script. The tree had changed under the
experiment.

\textbf{Case 11 (E81).} A static rule dictionary reported a real-to-shuffled
ratio of 0.96, read as ``finds no structure''. Decomposed by cue arity, the
two instruments inside it were opposite: multi-word cues 15:1,
single-word cues 0.92:1. The aggregate hid a working instrument and a
broken one.

\textbf{Case 12 (E10).} The pre-flight power check returned STOP. The run
proceeded anyway, because the author of the check overrode it. A control
its author ignores is not a control.

\section{Layer B: the control taxonomy}
\label{sec:taxonomy}

\subsection{What is a construction-matched null, and what is not}

Not every safeguard is a null, and calling them all ``construction-matched''
empties the term. We separate them:

\begin{itemize}
\item \textbf{Construction-matched null} (mechanism null): preserves the
  construction, perturbs the claimed mechanism. Cases 1, 2, 3, 9.
\item \textbf{Negative control}: an arm in which the treatment should have
  no effect, run to detect leakage or accidental signal. Case 4.
\item \textbf{Adversarial fixture}: input deliberately constructed to
  trigger a rule wrongly. Cases 2, 3.
\item \textbf{Implementation invariant}: an assertion about what the code
  must be doing, checked at runtime. Case 7.
\item \textbf{Measurement sanity check}: a check that the metric reads the
  field it claims to read. Cases 5, 8, 11.
\item \textbf{Replication check}: a redundant count or scan that a
  duplicate or drift would break. Case 6.
\item \textbf{In-run baseline}: a baseline recomputed in the same
  execution rather than cited. Case 10.
\item \textbf{Power gate}: a pre-flight check that the design can resolve
  the question. Case 12.
\end{itemize}

Only the first is a null in the strict sense. The others are safeguards,
and they are not interchangeable: an implementation invariant cannot detect
an invalid control, and a power gate cannot detect a parser reading the
wrong field.

\subsection{Detection status}

Not every safeguard fires at the same point, and it matters which:

\begin{itemize}
\item \textbf{Prevented before execution}: Case 12, partially --- the gate
  returned STOP and the run proceeded, so the gate prevented a verdict, not
  an execution. We do not count this as prevention.
\item \textbf{Detected during execution}: Cases 5, 6, 7, 8, 10.
\item \textbf{Detected after execution, before publication}: Cases 1, 2, 3,
  9, 11.
\item \textbf{Discovered retrospectively}: Case 4, where the control's
  validity was only established by a later model-scale comparison.
\end{itemize}

The distribution is itself a finding: five of twelve were caught only after
the number existed, which is why the workflow in Section~\ref{sec:workflow}
makes the null a required input rather than a post-hoc review.

\subsection{Claim language}

We make no claim that these errors are peculiar to AI-assisted research.
Circular benchmarks, measurement bugs, baseline drift, underpowered studies
and selective interpretation are all long documented in human research;
Section~\ref{sec:related} cites the classics.

The defensible difference is structural, not anthropological: AI-assisted
research compresses experiment design, implementation, execution and
interpretation into one machine-mediated loop, so such errors can be
generated and propagated rapidly and repeatedly within a single session.
That is a claim about the speed and coupling of the loop, and it is the
only one our evidence supports.

\section{Layer C: the executable workflow}
\label{sec:workflow}

\subsection{Structures}

\begin{verbatim}
Claim {
  id, hypothesis, predicted_effect, metric,
  falsification_condition, control_specification,
  data_hash, code_commit, timestamp, provenance, status
}
Experiment {
  id, claim_id, code_commit, data_hash,
  treatment, controls[], preregistered_decision_rule, outputs[]
}
Verdict = Supported | Contradicted | PartiallySupported
        | Underpowered | Artifact | Superseded
RevisionEdge { previous_claim, experiment, verdict, successor_claim }
\end{verbatim}

\subsection{The invariant}

\begin{quote}
\textbf{No result may modify its own decision rule after observation.}
\end{quote}

This is the invariant that separates pre-registration from
post-hoc rationalisation. We state its status honestly: the current system
\textbf{records} violations; it does not \textbf{prevent} them. Case 12 is
the proof --- an agent overrode the pre-flight gate and the run proceeded.
The correct description is therefore ``protocol violation recorded'', and
an earlier draft's claim that the workflow ``caught'' this case overstated
what happened.

This matters because a control system that is described as preventing
violations will be trusted more than one described as logging them, and the
difference is exactly what an agent under pressure will exploit.

\subsection{Failed branch retention}

The workflow's most distinctive property is that superseded claims and
failed experiments stay in the record. Figure~\ref{fig:lineage} shows a
real lineage from the ontology track, where every step is retained and the
revisions are visible.

\begin{figure}[h]
\centering
\begin{verbatim}
claim: geometry can certify ontology membership
  |
  +-- E12  anchor audit           -> refutes the fix it was designed to prescribe
  +-- E13  anchor ceiling         -> cancels E12's own recommendation
  +-- E14  injection realism      -> 0.880 becomes 0.506; no working certifier
  +-- E15  structure survives?    -> reopens the question: CONSTRAINT 0.612
  +-- E16  real misfilings        -> E15 withdrawn; everything at chance 0.519
  +-- E17  ANCHORVOC confound     -> three tests disagree; artifact
  +-- E18  ANCHORVOC closed       -> entry-length artifact; not a certifier
  |
  v
surviving claim: geometry is unreliable for certification
  |
  +-- E21  the question reversed  -> geometry as PROPOSER: top-3 0.668
  |
  v
surviving claim: unreliable judge, useful candidate generator
\end{verbatim}
\caption{The retained lineage. Every refuted intermediate survives in the
repository; the surviving claim is the bottom of the chain, not a rewrite
of the top.}
\label{fig:lineage}
\end{figure}

Reading the chain downward: a claim was made, eight experiments tested it,
seven failed, and each failure is preserved with its own result file. The
final claim is narrower than the first and was reachable only by keeping
the failures. This is the concrete form of ``failures stay visible''.

\subsection{What the workflow caught}

\begin{itemize}
\item \textbf{E14}: registered with a nearest-cell null; produced 0.506;
  marked Contradicted. Without the pre-registered null the 0.880 would have
  shipped.
\item \textbf{E83}: registered with a fixture it did not generate; produced
  recall 0.143; marked Artifact.
\item \textbf{E85}: registered with adversarial fixtures; produced precision
  0.000; marked Contradicted.
\item \textbf{E17}: registered with three tests and a rule that all must
  pass; test 3 failed; marked Artifact even though tests 1 and 2 passed.
\item \textbf{E10}: power gate returned STOP, run proceeded; recorded as a
  protocol violation; marked Underpowered.
\end{itemize}

In four of these five the outcome changed from positive to negative or
artifact. That is the workflow's value, and it is also its cost: it
converted apparent successes into documented failures.

\section{Falsifying the verifier}
\label{sec:falsify}

The apparatus this paper proposes is itself an artifact, so it gets the same
treatment it prescribes for everything else. This section reports the cases
where the verifier failed, because a checker with no kill tests is a
comfortable way of not finding things.

\subsection{The executable ledger, and attempts to violate it}

The workflow's invariants are executed, not documented. The current suite is
\textbf{nine invariants, zero failures}
(\texttt{research/methodology\_validation.json}, checked against 15 registered
claims):

\begin{enumerate}
\item \textbf{I1 threshold frozen} --- no result may change its own decision
  rule after observation.
\item \textbf{I2 null registered} --- every claim carries a
  construction-matched null before it runs.
\item \textbf{I3 STOP prevents execution} --- a power-gate STOP must not be
  followed by a run.
\item \textbf{I4 superseded addressable} --- a superseded claim must still
  resolve to something, not vanish.
\item \textbf{I5 artifact resolves} --- every claim's result file exists.
\item \textbf{I6 provenance} --- every claim records a source script and a
  commit.
\item \textbf{I7 manuscript numbers have IDs} --- no headline number lives
  only as hand-copied prose.
\item \textbf{I8 missing provenance fails} --- validation must refuse an
  incomplete record.
\item \textbf{I9 artifact parameters match} --- catches an artifact silently
  replaced by a differently-parameterised run.
\end{enumerate}

I8 is the one that proves the suite can fail. A deliberately stripped record
is submitted to it on every run, and it must reject it; it currently catches
three violations from that probe (missing artifact file, absent
\texttt{source\_script}, absent \texttt{git\_commit}). Without that probe,
``zero failures'' would be indistinguishable from a suite that never fires.

I3 is the invariant the workflow cannot satisfy. Case~12 overrode the
pre-flight gate and the run proceeded; the item above records the violation
rather than preventing it. The suite is honest about this: I3 reports
\texttt{executed: false} against the live ledger, not an enforcement count.

\subsection{Where the verifier itself was wrong}

Four failures, each found by using the thing rather than reviewing it:

\begin{enumerate}
\item \textbf{A stale artifact overwrote a corrected one.} A fixed output
  path plus a long-running process is a race: the process holds the code it
  started with, so a late-finishing run writes the \emph{old} verdict over
  the corrected artifact. This happened \textbf{four times in one session}
  on the Empedocles finite-size sweep, and once it printed the opposite
  conclusion. The fix was not a faster editor but an overwrite guard: the
  artifact now records \texttt{classifier\_version}, and a process running an
  older classifier refuses to replace a newer artifact, writing a side file
  instead.
\item \textbf{The interpreter contradicted its own numbers.} The first
  version of that same classifier emitted ``consistent with finite-size
  critical scaling'' while the diagnostics it had just computed showed the
  susceptibility peak sitting at the edge of the parameter sweep and a
  transition width that \emph{broadened} with system size. Both facts are
  incompatible with the verdict the classifier printed. The classification
  logic was rewritten (version 2, which tests the boundary condition
  explicitly and returns ``NO TRANSITION LOCATED''). This is the cleanest
  available example of why a machine-generated interpretation is a
  measurement like any other, and it is cross-referenced from Paper~10.
\item \textbf{The manuscript checker silently blessed a wrong number.} The
  script that traces Paper~4's figures back to result files contained a dead
  index expression (\texttt{[-1][-1] if False else ...}) that made one
  arithmetic derivation unreachable, so a number the paper had withdrawn
  (\texttt{+0.576}) was still being accepted as explained. The checker
  reported success. It was only found by asking where a specific printed
  number came from. Its replacement asserts the derivation instead of
  tolerating it.
\item \textbf{The new checker had two bugs of its own.} The cross-rater
  consistency checker was wrong on first run in two independent ways: it
  counted distinct entry names over the wrong row subset, and it built a
  rectangular contingency matrix for a chance-agreement term that requires a
  square one (the two raters use 57 and 42 categories respectively). Both
  were caught by running it, not by reading it, and both are now pinned by
  assertions.
\end{enumerate}

\subsection{Provenance that is present and wrong}

An invariant can pass while the record it checks is stale. The ledger
snapshot used above records commit \texttt{8232042}; the repository has
since moved on. Nothing in the suite fails on this, because the invariant
checks that provenance \emph{exists}, not that it is \emph{current}. That
gap is deliberate and named: a claim whose artifact predates the last
relevant code change should be re-run, and detecting that requires a
staleness relation the ledger does not yet carry.

\subsection{The principle}

A verifier needs its own kill tests, in the same sense that a classifier
needs a positive control. Concretely:

\begin{itemize}
\item every invariant must have a case that makes it fail, exercised on
  every run (I8's stripped-record probe);
\item every artifact path must be version-guarded against being overwritten
  by an older producer;
\item every derived number in a manuscript must re-derive, not merely
  reconcile to a tolerance;
\item every automated verdict must be compared against the numbers that
  produced it, because those are what the reader can check.
\end{itemize}

The Empedocles control battery (Paper~10) makes the same point from the other
side: an instrument that cannot fire is not a falsifier. This section is the
claim that the apparatus generating that control battery is subject to the
same rule.

\section{The empirical basis}

The catalog is not assembled by inspection from memory. Each case resolves
to a stored result file, and where a result file had been lost it was
regenerated before being relied on. Two of the E15--E18 experiments had
shipped as printing scripts with no saved output, so their numbers resolved
to nothing; both were re-run from the surviving corpus and ground truth and
reproduced their reported values (COSINE 0.882 under random injection and
0.519 on real misfilings; CONSTRAINT 0.612 and 0.536). The re-run is stored
in \texttt{research/e15\_e18\_rerun.json}.

Two claims in the original reports remain unreproducible because their
experiments have no saved output anywhere on disk (E5's 0.447 and E24's
0.712). They are not load-bearing here and we do not rely on them.

This is the workflow applied to the paper that describes the workflow: the
count of unreproducible claims went from five to two, and the two that
remain are named.

\section{Related work}
\label{sec:related}

\textbf{Exploratory versus confirmatory analysis.} The requirement to
separate hypothesis-generating from hypothesis-testing analysis is
long-established \cite{tukey1980we}, and the cost of undisclosed
analytic flexibility is well documented \cite{simmons2011false}. Our
contribution is not the principle but its automation: the null is a
required input to the execution rather than a review criterion applied
afterwards.

\textbf{Autonomous research systems.} The AI Scientist
\cite{lu2024aiscientist} implements an end-to-end loop that generates
hypotheses, runs experiments and writes up results. The Robot Scientist
\cite{king2004robot} automates hypothesis generation and experimentation
in a laboratory, with hypotheses formalised in logic. Both describe the
loop; neither describes the falsification infrastructure --- how the null is
specified before running, how a control's validity is itself checked, and
how failed branches are retained. Our contribution is that infrastructure.

\textbf{Belief revision and truth maintenance.} Our claim/verdict/revision
records are \emph{AGM-inspired} and \emph{related to} the belief-revision
framework of Alchourr\'on and Makinson \cite{agm1982logic}, and our
justification-and-retraction structure is related in spirit to Doyle's
truth maintenance system \cite{doyle1979tms}. We do not claim to implement
either. The AGM postulates are formal properties of contraction and
revision functions on belief sets; we have not shown that our revision
edges satisfy them, and our objects are empirical claims carrying effect
sizes and nulls, not logical propositions. The relationship is that of an
operational analogue, and we state it as such.

\textbf{Distribution fitting as a control.} One of our cases turns on
distinguishing a power law from a lognormal, for which we use standard
methods \cite{clauset2009powerlaw}. We mention it because a
heavy-tailed-looking sample is the most common route by which a null
result is mistaken for a positive one, and because a fit compared against
only one alternative family is not a test.

\section{Position against the current verification-gap literature}

A recent survey of autonomous research agents codes 26 systems across seven
audit dimensions and finds that in 24 runnable systems, \textbf{83\% release
code while only 38\% release the seeds or execution traces needed to reproduce
a run, and 38\% report any novelty-verification method}; among nine closed-loop
autonomous systems, seven are mechanical reruns and none demonstrates an
externally validated in-loop oracle \cite{ding2026verificationgap}. A companion
position paper argues that verification, attribution and reproducibility
infrastructure break down at agent scale \cite{mo2026paradigm}, and cites four
observed failure modes in AI-scientist systems: benchmark cherry-picking, data
leakage, metric misuse and post-hoc selection bias.

The present work is best read as a concrete instance of that gap and a partial
answer to it at the scale of one programme. The cases above are exactly the
failure modes that literature names --- leakage (the withdrawn Empedocles
cold-start curve), post-hoc selection (Case~12's override of the pre-flight
gate), and metric misuse (the classifier that reported critical scaling against
its own numbers) --- but each is recorded with an artifact, a verdict, and a
retained failed branch, which is the disclosure the survey reports as rare.

We do not claim to close the gap. Our ledger holds twelve cases from one
programme and is maintained by the same agent that runs the experiments. What
the survey's coding dimensions do not audit is whether a released verifier was
itself falsified; Section~\ref{sec:falsify} is an attempt to answer that
question for this one.

\section{Limitations}

\textbf{No control group.} We cannot claim the workflow produces better
research than research without it. That would require two groups using
different workflows on comparable problems, which we have not run. Every
statement here is about what the workflow recorded, not about comparative
outcomes.

\textbf{Agent compliance.} The workflow catches errors only if the agent
follows it, and Case 12 shows an agent can bypass it. It records the
violation; it does not prevent it.

\textbf{One programme, one domain.} All twelve cases come from one
research programme, one user, one machine fleet, over roughly three months.
The frequency of these failure modes elsewhere is unknown, and the
taxonomy's coverage of failure modes we did not happen to hit is unknown.

\textbf{Selection.} The cases are ones where a control \emph{caught}
something. We have no systematic study of cases where a control failed to
catch an error a human would have caught, nor of errors that survived to
publication.

\textbf{Detection-status coding.} The phase and timing labels in
Table~\ref{tab:catalog} are our own judgments, made retrospectively. They
are not independently coded.

\textbf{Novelty.} Each individual practice here is old. Pre-registration,
negative controls, power analysis, adversarial testing, in-run baselines
and retaining failed results are all standard, and a reviewer is right to
say so. We claim no invention of any of them. What we claim is narrower:
a precise definition of the construction-matched null that distinguishes
it from seven other safeguards commonly called ``controls''; twelve
documented cases from one programme with the correction chain intact; and
an honest accounting of which safeguards detected which errors, showing
that five of twelve were caught only after the number existed. If the
reader's question is ``is this new?'', the answer is no; if it is ``what
does it look like when an AI-driven loop actually goes wrong, and what
catches it?'', we have twelve answers with stored artifacts.

\section{Conclusion}

Twelve documented cases from one research programme show a consistent
shape: a strong positive result produced by a loop that also produced the
error, and a control that killed the positive. From those cases we derived
a precise definition of the construction-matched null, a taxonomy that
distinguishes a null from the seven other safeguards it is commonly
confused with, and an executable workflow whose records make the sequence
claim $\rightarrow$ experiment $\rightarrow$ verdict $\rightarrow$
revision auditable.

The workflow's most valuable property is the least glamorous: it retains
failed branches. The ontology lineage in Figure~\ref{fig:lineage} reaches
a defensible claim only because eight refuted intermediates were kept. If
they had been deleted --- or never written down --- the surviving claim
would have looked like a first guess rather than the end of a search.

We state plainly what this is not: not a proof that AI-assisted research is
worse than human research, not an implementation of AGM belief revision,
and not a prevention system. It is an instrumentation proposal with twelve
worked examples, and its main result is that in four of the five cases where
it was applied end-to-end the apparent positive conclusion was overturned,
reclassified as an artifact, or rendered non-promotable --- by a control, a
pre-registered null, a provenance failure or a protocol violation, and never
by a change of preference. We avoid the phrase ``the answer changed sign''
because the outcomes are not one kind of event: an artifact, an underpowered
run, a protocol violation and a reclassification are different failures and
conflating them would hide which safeguard did the work.

\begin{thebibliography}{99}

\bibitem[Ding et al.(2026)]{ding2026verificationgap}
Ding, T., Nannapaneni, A., Liu, B., \& Zhang, L. (2026).
\newblock Autonomous research agents: a survey of AI scientists and the
verification gap.
\newblock arXiv:2608.05179.

\bibitem[Mo(2026)]{mo2026paradigm}
Mo, B. (2026).
\newblock Position: the age of AI agents demands a new scientific paradigm to
sustain trustworthy science.
\newblock arXiv:2607.26064.

\bibitem[Alchourr\'on \& Makinson(1982)]{agm1982logic}
Alchourr\'on, C.~E., \& Makinson, D.~C. (1982).
\newblock On the logic of theory change: Contraction functions and their
associated revision functions.
\newblock \emph{Theoria}, 48(1), 14--37.
\newblock DOI: 10.1111/j.1755-2567.1982.tb00480.x.

\bibitem[Clauset et al.(2009)]{clauset2009powerlaw}
Clauset, A., Shalizi, C.~R., \& Newman, M.~E.~J. (2009).
\newblock Power-law distributions in empirical data.
\newblock \emph{SIAM Review}, 51(4), 661--703.
\newblock DOI: 10.1137/070710111.

\bibitem[Doyle(1979)]{doyle1979tms}
Doyle, J. (1979).
\newblock A truth maintenance system.
\newblock \emph{Artificial Intelligence}, 12(3), 231--272.
\newblock DOI: 10.1016/0004-3702(79)90008-0.

\bibitem[King et al.(2004)]{king2004robot}
King, R.~D., Whelan, K.~E., Jones, F.~M., Reiser, P.~G.~K., Bryant, C.~H.,
Muggleton, S.~H., Kell, D.~B., \& Oliver, S.~G. (2004).
\newblock Functional genomic hypothesis generation and experimentation by a
robot scientist.
\newblock \emph{Nature}, 427(6971), 247--252.
\newblock DOI: 10.1038/nature02236.

\bibitem[Lu et al.(2024)]{lu2024aiscientist}
Lu, C., Lu, C., Lange, R.~T., Foerster, J., Clune, J., \& Ha, D. (2024).
\newblock The AI Scientist: Towards fully automated open-ended scientific
discovery.
\newblock \emph{arXiv preprint} arXiv:2408.06292.

\bibitem[Simmons et al.(2011)]{simmons2011false}
Simmons, J.~P., Nelson, L.~D., \& Simonsohn, U. (2011).
\newblock False-positive psychology: Undisclosed flexibility in data
collection and analysis allows presenting anything as significant.
\newblock \emph{Psychological Science}, 22(11), 1359--1366.
\newblock DOI: 10.1177/0956797611417632.

\bibitem[Tukey(1980)]{tukey1980we}
Tukey, J.~W. (1980).
\newblock We need both exploratory and confirmatory.
\newblock \emph{The American Statistician}, 34(1), 23--25.
\newblock DOI: 10.1080/00031305.1980.10482706.

\end{thebibliography}

\end{document}
